%%%%%%%%%%  scan idx 307  |  folio 294  %%%%%%%%%%
\textbf{3.3. Groupes de Grothendieck relatifs}

\textbf{3.3.1.} Soit \(f:X\longrightarrow Y\) un morphisme localement de type fini de schémas. Nous poserons~:
\[
\begin{gathered}
 K^{\bullet}(f)=k(D(f)_{\parf}),\\
 K_{\bullet}(f)=k(D^b(f)_{\coh}),
\end{gathered}
\]
où \(D^b(f)_{\coh}\) désigne la catégorie des complexes pseudo-cohérents rel. à \(f\) et à cohomologie bornée (III 1.2), et \(D(f)_{\parf}\) la catégorie des complexes d'amplitude parfaite finie rel. à \(f\) (III 4.2). Quand il n'y aura pas d'ambiguïté possible, nous écrirons parfois \(K^{\bullet}(X/Y)\) (resp. \(K_{\bullet}(X/Y)\)) au lieu de \(K^{\bullet}(f)\) (resp. \(K_{\bullet}(f)\)).

Les groupes \(K^{\bullet}(f)\) et \(K_{\bullet}(f)\) ne possèdent pas la riche structure de leurs analogues «~absolus~» \(K^{\bullet}(X)\) et \(K_{\bullet}(X)\), car le produit tensoriel ne définit pas en général d'homomorphisme \(K^{\bullet}(f)\otimes K^{\bullet}(f)\longrightarrow K^{\bullet}(f)\) (resp. \(K^{\bullet}(f)\otimes K_{\bullet}(f)\longrightarrow K_{\bullet}(f)\)). En revanche, ils vérifient d'intéressantes propriétés de variance, que nous allons passer en revue.

\textbf{3.3.2.} Soit
\[
\begin{tikzcd}
X\arrow[rr,"f"]\arrow[dr,"h"'] &&Y\arrow[dl,"g"]\\
&S&
\end{tikzcd}
\]
un triangle commutatif de morphisme localement de type fini de schémas.
\begin{enumerate}[label=\textbf{3.3.2.\arabic*.}, leftmargin=4em]
\item Si \(g\) est pseudo-cohérent, on a \(K_{\bullet}(f)=K_{\bullet}(h)\) (III 1.12). Si \(g\) est lisse et \(X\) quasi-compact, on a \(K^{\bullet}(f)=K^{\bullet}(h)\) (III 4.4.1).

\item D'après (III 1.11 et 4.5), le foncteur \(\otimes_Y^{\mathrm L}\) définit des accouplements
\[
 D(f)_{\parf}\times D^b(g)_{\coh}\longrightarrow D^b(h)_{\coh},
\]
\[
 D(f)_{\parf}\times D(g)_{\parf}\longrightarrow D(h)_{\parf},
\]
qui induisent des homomorphismes
\[
 K^{\bullet}(f)\otimes K_{\bullet}(g)\longrightarrow K_{\bullet}(h),
\]
\[
 K^{\bullet}(f)\otimes K^{\bullet}(g)\longrightarrow K^{\bullet}(h).
\]
